
% ==========================================================
% Study H27: From axioms to lemmas + strengthened bridge
% ==========================================================

\section{Study H27: Von Axiomen zu Lemmas (Spin-Zwang), und eine st\"arkere Br\"ucke Richtung RH}

\subsection{Zielsetzung}
Wir ersetzen die in H26 verwendeten \emph{Axiome} schrittweise durch \emph{ableitbare Lemmas}.
Das Ziel ist eine BM-interne Kette
\[
\text{Off-Axis} \ \Rightarrow\ \text{wachsender Twist} \ \Rightarrow\ \text{Morley-Verletzung} \ \Rightarrow\ \bot,
\]
und parallel dazu eine Formulierung der Zwei-Komponenten-Hohlraumidee als \emph{Gate-Zerlegung} des Planck-Spektrums.

\subsection{(I) Geometrische Bausteine: Kontraktion und Drehung (2er/3er)}
\paragraph{2er-Kontraktion (``Halbierung'').}
Wir definieren einen Operator $\mathcal C_2$ auf einem (orientierten) BM-Tetraeder $\Delta$:
\[
\mathcal C_2(\Delta;R) := \text{``inskribiere einen inneren Tetraeder mit Radius }R/2\text{''}.
\]
Die zugeh\"orige Drehung (Spinfixierung) sei eine Rotation um $\pi$ um eine ausgezeichnete Achse:
\[
\mathcal R_\pi: \Delta \mapsto \Delta^\pi.
\]
Die kombinierte 2er-Regel lautet:
\[
\Delta \xrightarrow{\ \mathcal C_2\ }\Delta_{1/2}\xrightarrow{\ \mathcal R_\pi\ }\Delta_{1/2}^\pi.
\]
\emph{Interpretation:} Die in den Rechnungen auftauchende ``eineinhalb''-Achse kann als \emph{Zweier-Kontraktion plus Drehung}
verstanden werden: nicht nur L\"angenhalbierung, sondern \emph{Achsenneubildung} durch den $\pi$-Twist.

\paragraph{3er-Kontraktion (``Drittelung'' + $180^\circ$-Fixierung).}
Analog definieren wir $\mathcal C_3$:
\[
\mathcal C_3(\Delta;R):=\text{``innerer Tetraeder mit Radius }R/3\text{''},\qquad
\mathcal R_\pi: \Delta \mapsto \Delta^\pi.
\]
Die beobachtete ``ABC-Elementarit\"at'' wird als Stabilit\"atsbedingung formuliert:
\[
\Delta \xrightarrow{\ \mathcal C_3\ }\Delta_{1/3}\xrightarrow{\ \mathcal R_\pi\ }\Delta_{1/3}^\pi,
\quad \text{wobei ABC als Dreierstruktur invariant bleibt.}
\]

\paragraph{Glatte Zahlen / Hamming-Zahlen.}
Die Halbr\"aume erzeugt durch wiederholte Anwendung von $\mathcal C_2,\mathcal C_3$ eine Skalenhalbgruppe
\[
\langle 2^{-a}3^{-b} : a,b\in\mathbb N_0\rangle,
\]
was genau die glatten Skalen (Hamming) repr\"asentiert.
BM-intern wird damit die Entstehung der ``2'' und ``3'' als Skalenoperatoren formalisiert.

\subsection{(II) Oktonischer Spin-Zwang als Winkelhalbierenden-Lemma}
\paragraph{Spin-Zwang (informell).}
Im oktonischen Erweiterungsschritt ist die Multiplikation nicht voll assoziativ; zul\"assig bleiben nur Produkte, deren
Resultat geometrisch auf einer \emph{Winkelhalbierenden} liegt (``erzwungener Spin'').

\begin{lemma}[Winkelhalbierenden-Kriterium (BM-Spin)]
Seien $a,b,c$ die drei elementaren Achsen (ABC) im BM. Die Operationen $a\cdot b\mapsto c$, $a\cdot c\mapsto b$, etc.
sind als \emph{zul\"assige additive Updates} genau dann stabil, wenn die entstehende Phase
auf einer oktonischen Winkelhalbierenden liegt. Abweichungen erzeugen eine Phasendrift $\Delta\theta$,
die sich additiv entlang der Skala akkumuliert.
\end{lemma}

\paragraph{Kosinus-$120^\circ$ und L\"angenhalbierung.}
Die beobachtete Halbierung kann aus dem Vektoradditionsgesetz bei $120^\circ$ motiviert werden:
\[
\|u+v\|^2=\|u\|^2+\|v\|^2+2\|u\|\|v\|\cos 120^\circ
=\|u\|^2+\|v\|^2-\|u\|\|v\|.
\]
F\"ur passende Normierung ergibt dies effektive Kontraktionsfaktoren, die als ``neue Achse'' wirken.

\subsection{(III) Morley-Abweichung ist Lipschitz in den Winkeln (Stabilit\"atslemma)}
Wir betrachten die Einbettung der BM-Phasen als Punkte auf dem Einheitskreis:
\[
P(\theta)=(\cos\theta,\sin\theta).
\]
Seien $A=P(\theta_A)$, $B=P(\theta_B)$, $C=P(\theta_C)$.

\begin{lemma}[Lipschitz-Stabilit\"at der Morley-Abweichung]
F\"ur nichtdegenerierte Dreiecke (Winkelabst\"ande $\ge \eta>0$) existiert eine Konstante $L(\eta)>0$ so dass
\[
|\delta_M(A,B,C(\theta_C+\Delta\theta))-\delta_M(A,B,C(\theta_C))|
\ \ge\ L(\eta)\,|\Delta\theta|
\]
f\"ur alle hinreichend kleinen $\Delta\theta$ gilt.
\end{lemma}

\paragraph{Konsequenz.}
Wenn ein Off-Axis-Mechanismus einen Twist $\Delta\theta$ erzeugt, ist die Morley-Verletzung im BM \emph{mindestens linear}
in $|\Delta\theta|$, solange man Dreiecke mit guter Konditionierung w\"ahlt.

\subsection{(IV) Existenz gut konditionierter Dreiecke (adversarial/Probabilistik)}
\begin{lemma}[Existenz eines ``Zeugen''-Dreiecks]
Unter einer milden Gleichverteilungsannahme der BM-Phasen in einer Zelle gibt es f\"ur jedes $\eta\in(0,\pi/3)$
und hinreichend gro\ss es Indexset $\mathcal I_{m,\lambda}$ ein Tripel $(A,B,C)$ mit Winkelabst\"anden $\ge \eta$.
\end{lemma}

\paragraph{Kommentar.}
Das ist die formale Version dessen, was wir in H23/H24 numerisch als adversarial Suche umgesetzt haben:
statt \emph{ein beliebiges} Tripel zu nehmen, w\"ahlt man ein gut konditioniertes Tripel, das Twist-Effekte stark
in $\Delta\delta_M$ \"ubersetzt.

\subsection{(V) Cutoff-Lemma aus Twist-Wachstum}
In H24 wurde Off-Axis als wachsender Twist modelliert:
\[
\Delta\theta(\gamma)=\varepsilon\,\frac{\gamma}{\gamma_0}\pi.
\]
Mit den Lemmas aus (III)--(IV) folgt:

\begin{lemma}[Cutoff-Existenz aus wachsendem Twist]
Es existiert ein $T^\*$, so dass f\"ur $T\ge T^\*$ in einem positiven Anteil der Zellen
\[
\Delta\delta_M^{\max}(T)>\tau
\]
gilt. Insbesondere w\"achst $\Delta\delta_M^{\max}(T)$ mindestens linear in $T$ bis zur Schranke $\tau$.
\end{lemma}

\paragraph{Numerischer Anker (H24).}
Mit $\tau=2\hat h\approx 0.01818$ ergab sich:
\[
T_{50\%}\approx 800,\qquad T_{90\%}\approx 1400.
\]

\subsection{(VI) St\"arkere Br\"ucke als reines Index-Matching: Spektral-Residual aus Primdaten}
Der Index-Matching-Schritt $p^\*(T)=p_{N(T)}$ ist praktisch, aber als Br\"ucke schwach.
Eine st\"arkere Formulierung ist: definiere eine \emph{prime-Observable} $\mathcal O(X)$ aus den Primzahlen, die im BM
direkt mit dem Spektralterm gekoppelt ist.

\paragraph{Vorschlag (Explizite-Formel-Br\"ucke).}
Nutze eine Chebyshev/``prime power''-Observable, z.\,B.
\[
\psi(X):=\sum_{p^k\le X}\log p.
\]
Die klassische explizite Formel (schematisch) lautet
\[
\psi(X)=X-\sum_{\rho}\frac{X^{\rho}}{\rho}+\text{(bekannte Neben\-terme)}.
\]
Ein Off-Axis-Nullstelle $\rho=\beta+i\gamma$ mit $\beta>\tfrac12$ erzeugt einen Beitrag der Gr\"o\ss enordnung
\[
\left|\frac{X^\rho}{\rho}\right|\sim \frac{X^\beta}{|\rho|},
\]
der f\"ur gro\ss es $X$ schneller w\"achst als jeder $\beta=\tfrac12$-Beitrag.

\begin{lemma}[BM-Residual aus $\psi(X)$]
Definiere den BM-Residual
\[
\mathcal R(X):=\frac{|\psi(X)-X|}{X^{1/2}}.
\]
Dann bedeutet Off-Axis ($\beta>\tfrac12$) heuristisch:
\[
\mathcal R(X)=\Omega\!\bigl(X^{\beta-1/2}\bigr),
\]
also unbeschr\"ankt wachsend.
Wenn BM hingegen (Spin-Zwang) eine Schranke $\mathcal R(X)\le C(\tau)$ f\"ur alle $X$ fordert,
entsteht ein Widerspruch \emph{ohne} Prime-Cutoff.
\end{lemma}

\paragraph{Interpretation.}
Damit ersetzt man ``Prime-Cutoff'' durch eine klassische, arithmetische Form: Off-Axis w\"urde eine
\emph{zu gro\ss e} Abweichung in einer Prime-Observable erzwingen.

\subsection{(VII) Zwei-Komponenten-Planck: Gate und Datenindikator}
Wir hatten in H26 die Gate-Zerlegung
\[
u_{\mathrm{b}}=u_P\cdot g,\qquad u_{\mathrm{d}}=u_P\cdot(1-g),
\]
mit $g$ als sigmoidale Stufe und $w\to 0$ als ``y-Achsen-Parallele''.
Datenindikator: in H20 zeigt die \emph{dark fraction} einen Regimewechsel bei
\[
\lambda^\*\approx 0.20.
\]
Zus\"atz: die Click-Rate (H13) zeigt entlang $\lambda$ ebenfalls Struktur, allerdings nur moderat konsistent
(auf 6 St\"utzstellen ist dies statistisch zart; Korrelationen sind klein).

\subsection{Konkrete n\"achste Beweiselemente (H28)}
Damit das Schema wirklich \emph{beweisbar} wird, sind drei harte Schritte zu liefern:
\begin{enumerate}
\item \textbf{Geometrie $\Rightarrow$ Morley-Schranke:} eine explizite Absch\"atzung $L(\eta)$ und ein Beweis, dass BM-Dreiecke
in den relevanten Zellen mit $\eta$-Konditionierung existieren.
\item \textbf{Off-Axis $\Rightarrow$ Twist-Wachstum:} Ableitung des Wachstums aus der oktonischen Nichtassoziativit\"at plus
2er/3er-Kontraktionslogik (``1/2 + Drehung'' als Achsenneubildung).
\item \textbf{Prime-Observable-Kopplung:} Festlegung einer BM-Observable $\mathcal O$ (z.\,B.\ aus $\psi(X)$ oder aus
Sch\"utte-/Click-Gating), die eine \emph{globale Schranke} besitzt, die Off-Axis verletzen w\"urde.
\end{enumerate}
